\section{Follow-up Question Ranking and Conversational Retrieval}
\label{sec:followup}

A causal embedding model for procedural documentation is not consumed in
isolation: in the AEP/AJO assistant it sits behind a conversational surface
that must decide \emph{what the user should ask next}. This section surveys the
research landscape underpinning that capability---follow-up reranking,
clarifying-question generation and ranking, multi-turn conversational
retrieval, and proactive/mixed-initiative dialogue---and argues that ordering
follow-ups by \emph{causal/next-step progression} rather than by mere semantic
relevance is the distinguishing requirement for our setting.

\subsection{The follow-up reranking task in this project}

The concrete task instance in the repository is a \emph{slate-reranking}
problem. The production data file
(\texttt{prod\_jan\_feb\_26-aep-ajo\_follow-up-queries.json}) is a list of
$334$ records, each of the form
\begin{quote}\small
\texttt{\{"question": <user query>,}\\
\texttt{ "follow\_up\_queries": [[ fq$_1$, \dots, fq$_{10}$ ]]\}}
\end{quote}
i.e.\ a single user question paired with a slate of exactly ten candidate
follow-up questions (the inner list is wrapped once). For the anchor
\emph{``What is Journey Pause feature''} the candidates range over permissions,
the hold-vs-discard semantics at action nodes, scheduled audience runs during a
pause, organization-wide held-profile limits, and re-entrance/timeout
interaction---a mix of \emph{prerequisite}, \emph{consequence}, and
\emph{lateral} questions. The reranking harness
(\texttt{rerank\_followup.py}) flattens every record into
\texttt{(question, follow\_up)} pairs, scores each pair with the fine-tuned
\texttt{DirectionalClassifier} bi-encoder
($\mathrm{score}=\sigma(\text{MLP}([a;b;a-b;a\!\cdot\!b]))$, the architecture
from Section~\ref{sec:embeddings}), and sorts the slate descending by score.
Crucially the model used is the \emph{directional precedence} classifier, so
the score approximates $P(\text{question causally precedes/enables follow-up})$
rather than topical similarity---this is exactly the ``rank by logical
progression'' signal we want, and it is what differentiates the task from a
generic relevance reranker. The slate is fixed and there is no gold ordering in
the file, so evaluation must come from a downstream signal (click logs,
LLM-as-judge ordering, or a held-out human ranking).

\subsection{Clarifying-question generation and ranking}

When a user query is under-specified, a proactive system can either retrieve or
\emph{ask}. \citet{fu:rao2018clarification} framed clarifying-question
\emph{ranking} as maximizing the expected value of perfect information (EVPI)
over StackExchange post--comment data, selecting the question whose likely
answer most changes the system's belief.
\citet{fu:aliannejadi2019clarifying} introduced the offline
\textbf{Qulac} methodology and the influential SIGIR'19 formulation of
\emph{asking clarifying questions in open-domain information-seeking
conversations}, later scaled into the \textbf{ClariQ}/ConvAI3 challenge
\citep{fu:aliannejadi2021convai3} with explicit \emph{clarification-need}
labels (the \texttt{clariq} TSV in our repo carries
\texttt{topic\_id, initial\_request, facet\_desc, question, answer} columns and
a 4-level clarification-need score).
\textbf{ClarQ} \citep{fu:kumar2020clarq} bootstrapped $\sim$2M
clarification questions across 173 StackExchange domains for generation, and
\textbf{CLAMBER} \citep{fu:clamber2024} contributes an LLM-oriented taxonomy of
$\sim$12K ambiguous queries, finding that current LLMs are still weak at both
\emph{recognizing} ambiguity and \emph{generating} a useful clarifying question.
At web scale, \textbf{MIMICS} \citep{fu:zamani2020mimics} mined real search
clarification panes (the ``People also ask''/clarification-pane lineage). These
threads share our core sub-problem: given a query, produce and \emph{rank} a
small slate of next questions; the AEP follow-up slate is the
documentation-grounded, ten-candidate version of this.

\subsection{Conversational and multi-turn retrieval}

Follow-ups live inside multi-turn sessions, where the central difficulty is
\emph{context-dependence}: a follow-up like ``what permissions does it need?''
is meaningless without its antecedent. Two solution families dominate.
\paragraph{Query rewriting.} \textbf{CANARD} \citep{fu:elgohary2019canard}
reformulates context-dependent questions into self-contained ones; T5-based
rewriters (\textbf{T5QR}, \citealp{fu:lin2020t5qr}) became the de-facto
sequence-to-sequence approach, turning a conversation into a standalone query
that any ad-hoc retriever can consume. \textbf{QReCC}
\citep{fu:anantha2021qrecc} packages 14K conversations as an end-to-end
``rewrite$\rightarrow$retrieve$\rightarrow$answer'' benchmark
(\texttt{qrecc\_train.json}/\texttt{qrecc\_test.json} in the repo).
\paragraph{Conversational dense retrieval.} Rather than rewrite, \textbf{ConvDR}
\citep{fu:yu2021convdr} trains a session encoder to embed the full multi-turn
context directly, distilling from the embedding of an oracle rewrite
(teacher--student) so it works few-shot. \textbf{TopiOCQA}
\citep{fu:adlakha2022topiocqa} stresses \emph{topic switching}: $3{,}920$
conversations averaging $13$ turns over $\sim$4 documents each, forcing
retrieval that detects when the user pivots---directly analogous to a user
jumping from ``pause a journey'' to ``read-audience scheduling.'' Most
recently, \textbf{MTRAG} \citep{fu:katsis2025mtrag} is a human-generated
multi-turn RAG benchmark ($110$ conversations, $\sim$7.7 turns, four domains
incl.\ Cloud) explicitly probing non-standalone and unanswerable later-turn
questions; the repo's \texttt{mtrag/conversations.json} mirrors this format with
per-turn retriever traces.

\subsection{Proactive and mixed-initiative agents}

Reranking a fixed slate is a special case of \emph{proactivity}: deciding,
unprompted, what to surface next. \citet{fu:deng2023mixedinitiative} survey
proactive dialogue (clarification, asking, recommending, non-collaborative
initiative). \textbf{ProactiveAgent} \citep{fu:lu2024proactive} (the
\texttt{proactive\_agent} dir) trains an agent on $\sim$6.8K real events to
\emph{anticipate} the user's next task and offer assistance before being asked,
with a reward model scoring whether an intervention was wanted. On the
generation side, \textbf{FollowupQG} \citep{fu:meng2023followupqg} (the
\texttt{followupqg} split: \texttt{\{question, answer, follow-up, relation\}}
tuples from Reddit ELI5) targets \emph{information-seeking} follow-up
generation, showing that model follow-ups remain markedly less informative than
human ones---evidence that good follow-ups require genuine reasoning about what
new knowledge the answer leaves open, not surface paraphrase.

\subsection{Reranking methods and the causal/next-step signal}

The reranker toolkit spans pointwise cross-encoders
(\textbf{monoBERT}, \citealp{fu:nogueira2019bertrerank}; \textbf{monoT5},
\citealp{fu:nogueira2020monot5}) that jointly attend over (query, candidate),
and listwise LLM rerankers (\textbf{RankGPT}, \citealp{fu:sun2023rankgpt};
distilled into the 7B \textbf{RankZephyr}, \citealp{fu:pradeep2023rankzephyr})
that emit a full permutation ``[3] > [1] > \dots'' over the slate. All of these
optimize \emph{relevance}. The thesis of this survey is that a follow-up slate
should instead be ordered by \emph{logical/causal progression}: the next
question a user \emph{should} ask is typically the immediate prerequisite or the
immediate consequence of the current step, not merely the most topically
similar one. Two ingredients combine cleanly: (i) the joint cross-attention of a
cross-encoder/LLM reranker, which Section~\ref{sec:embeddings} and the repo's
ordering experiments show generalizes better to non-adjacent pairs; and (ii) a
\emph{directional precedence} score $s(q,fq)\approx P(q \prec fq)$ from the
causal embedding model. A listwise reranker can be supervised or prompted to
sort by $s$ rather than by relevance, and the same pairwise scores feed the rank
aggregator of Section~\ref{sec:procedural} when a globally consistent ordering
of the slate is required.

\paragraph{Relevance to causal embedding for AEP/AJO.}
The follow-up slate is the most direct deployment of the causal embedding
primitive: \texttt{rerank\_followup.py} already scores
$(question, follow\_up)$ with the directional bi-encoder and sorts by it. The
literature surveyed here motivates three concrete upgrades for our
$<\!1$B-parameter model. First, treat the slate as a \emph{partial order} of
prerequisite/consequence/lateral follow-ups (mirroring CLAMBER's ambiguity
types and proScript's partial-order scripts in Section~\ref{sec:procedural})
rather than a single relevance ranking, so parallel ``lateral'' follow-ups are
not spuriously forced above causal next-steps. Second, adopt a cross-encoder or
distilled listwise reranker (RankZephyr-style) so the ordering signal benefits
from joint attention---consistent with the repo finding that the DeBERTa
cross-encoder dominates on \emph{ordering} despite trailing on flat
classification. Third, because follow-ups arrive mid-conversation, a
ConvDR/CANARD-style rewrite step should canonicalize the anchor question before
causal scoring, so context-dependent follow-ups (``how does \emph{it} interact
with re-entrance?'') resolve to their AEP/AJO referents before precedence is
estimated.
